\chapter{Future Work: Apple Silicon GPU Support}

\section{Current State}

This package currently supports Apple Silicon \textbf{CPUs} via the ARM/AArch64 backend:

\begin{lstlisting}[style=ts]
import { ARMCodegen } from "@euriklis/llvm-ir";

// Target Apple M3/M5 CPU cores
const codegen = new ARMCodegen("apple_module", true);  // 64-bit
codegen.targetTriple = "arm64-apple-macosx";

// Use NEON SIMD for vectorization (128-bit vectors)
const vec4 = new VectorType({ elementType: float, count: 4 });
// Generates efficient ARM NEON instructions
\end{lstlisting}

This works well for CPU-bound computations with SIMD acceleration.

\section{The Apple GPU Challenge}

Apple M3/M5 chips have powerful GPUs that are excellent for AI and parallel computing. However, unlike NVIDIA and AMD, Apple uses a \textbf{proprietary GPU architecture}:

\textbf{How GPU compilation works:}

\begin{itemize}[noitemsep]
  \item \textbf{NVIDIA}: Your code → LLVM IR (NVPTX) → PTX → CUDA binary (supported)
  \item \textbf{AMD}: Your code → LLVM IR (AMDGPU) → GCN/RDNA binary (supported)
  \item \textbf{Apple}: Your code → Metal Shading Language\cite{metalspec} → \textbf{AIR (proprietary)} → Metal binary (not supported)
\end{itemize}

\textbf{The problem:} AIR (Apple Intermediate Representation) is based on LLVM IR, but Apple keeps the format \textbf{closed and undocumented}\cite{airformat}. We cannot generate it directly.

\section{Why This Matters for TypeScript Developers}

If you're coming from a JavaScript/TypeScript background, here's an analogy:

\begin{itemize}[noitemsep]
  \item CUDA/AMD are like \textbf{open APIs}—you can generate code programmatically (what this package does)
  \item Metal is like a \textbf{proprietary framework}—you must use Apple's tools and pre-built libraries
\end{itemize}

Think of it like this:
\begin{lstlisting}[style=ts]
// CUDA/AMD (what this package does)
const kernel = generateLLVMIR({
  instructions: [...],  // Full programmatic control
  intrinsics: [...]
});

// Metal (what Apple requires)
const kernel = `
  kernel void myKernel(...) {
    // Must write Metal Shading Language as text
  }
`;
// Then compile with Apple's tools
\end{lstlisting}

\section{How PyTorch and TensorFlow Do It}

Frameworks like PyTorch don't generate Metal code either. Instead, they use \textbf{MPS (Metal Performance Shaders)}\cite{pytorchmps}\cite{applepytorch}—Apple's high-level library:

\begin{itemize}[noitemsep]
  \item Apple provides hundreds of \textbf{pre-optimized GPU kernels} (matrix multiply, convolution, etc.)\cite{metaldocs}
  \item PyTorch maps operations to these kernels via the \textbf{MPSGraph API}
  \item This works great for standard ML operations
  \item For custom kernels, you must write \texttt{.metal} files by hand\cite{metalspec}
\end{itemize}

\section{What We're Waiting For}

Apple has two paths to enable this package's Metal support:

\textbf{Option 1: Open the AIR format specification}

If Apple documents how AIR works\cite{airformat}, we could:
\begin{lstlisting}[style=ts]
import { MetalCodegen } from "@euriklis/llvm-ir";

const codegen = new MetalCodegen("my_kernel");
codegen.targetTriple = "air64-apple-macosx";
// Generate AIR directly, just like we do for NVPTX
\end{lstlisting}

\textbf{Option 2: Stable Metal Shading Language generation}

Alternatively, we could generate \texttt{.metal} source files programmatically\cite{metaltools}:
\begin{lstlisting}[style=ts]
const codegen = new MetalCodegen("my_kernel");
const metalSource = codegen.toMetalSource();  // Generate .metal text
// User compiles: xcrun metal -c kernel.metal
\end{lstlisting}

This is less elegant than LLVM IR generation, but would work.

\section{Current Workarounds}

Until Apple opens AIR or provides better tooling, TypeScript/JavaScript developers targeting Apple GPUs should:

\begin{enumerate}[noitemsep]
  \item \textbf{Use high-level frameworks}: PyTorch MPS, TensorFlow, Core ML
  \item \textbf{Write Metal Shading Language directly}: Create \texttt{.metal} files for custom kernels
  \item \textbf{Use this package for CPU SIMD}: ARM/NEON support works today
  \item \textbf{Prototype on NVIDIA/AMD}: Develop with this package's GPU support, deploy to Apple via Metal ports
\end{enumerate}

\section{Unified Memory Advantage}

One unique strength of Apple Silicon: \textbf{unified memory architecture}.

\begin{itemize}[noitemsep]
  \item CPU and GPU share the same memory pool (16-128 GB)
  \item No expensive CPU$\leftrightarrow$GPU copies (unlike NVIDIA)
  \item Ideal for workflows mixing CPU and GPU work
\end{itemize}

Once Metal support is added, this will make Apple Silicon extremely competitive for heterogeneous computing workflows.

\section{Timeline}

\textbf{When AIR becomes public or Apple provides stable LLVM IR → Metal tooling}, we will add:

\begin{itemize}[noitemsep]
  \item \texttt{MetalCodegen} class matching the API of \texttt{NVVMCodegen}, \texttt{AMDGPUCodegen}
  \item Metal-specific intrinsics (\texttt{threadgroup} memory, barriers, SIMD-groups)
  \item Complete examples for Apple GPU programming
  \item Validation with Apple's Metal compiler
\end{itemize}

We're monitoring Apple's developer documentation and LLVM commits for any changes. If you have insights into Metal's IR generation, please contribute!

\section{Community Contributions Welcome}

Possible paths forward:
\begin{itemize}[noitemsep]
  \item \textbf{Metal Shading Language generator}: Generate \texttt{.metal} source text (similar to how we generate LLVM IR text)
  \item \textbf{MPSGraph bindings}: High-level TypeScript API for Metal Performance Shaders
  \item \textbf{AIR reverse-engineering}: Document the format (challenging, legal concerns)
  \item \textbf{Lobby Apple}: Request official AIR documentation or LLVM Metal backend
\end{itemize}

If you're interested in helping, see the repository's contribution guidelines.

\backmatter

\begin{thebibliography}{99}
\bibitem{llvmessentials}
Suyog Sarda, Mayur Pandey. \textit{LLVM Essentials}. Packt Publishing, 2015. An excellent overview of LLVM fundamentals; this handbook extends those concepts to the TypeScript ecosystem and multi-architecture workflows.

\bibitem{llvmbindings}
ApsarasX. \textit{llvm-bindings: LLVM bindings for Node.js/JavaScript/TypeScript}. GitHub, 2021. Available at: \url{https://github.com/ApsarasX/llvm-bindings}

\bibitem{llvmnode}
Maël Nison et al. \textit{llvm-node: Node.js bindings for LLVM}. npm package, 2023. Available at: \url{https://www.npmjs.com/package/llvm-node}

\bibitem{nodellvmc}
Cornell Capra. \textit{node-llvmc: JavaScript/TypeScript FFI bindings for the LLVM C API}. GitHub, 2023. Available at: \url{https://github.com/cucapra/node-llvmc}

\bibitem{metalspec}
Apple Inc. \textit{Metal Shading Language Specification, Version 4}. Apple Developer Documentation, 2025. Available at: \url{https://developer.apple.com/metal/Metal-Shading-Language-Specification.pdf}

\bibitem{metaldocs}
Apple Inc. \textit{Metal Documentation}. Apple Developer, 2024. Available at: \url{https://developer.apple.com/metal/}

\bibitem{pytorchmps}
PyTorch Contributors. \textit{MPS backend --- PyTorch Documentation}. PyTorch, 2024. Available at: \url{https://docs.pytorch.org/docs/stable/notes/mps.html}

\bibitem{applepytorch}
Apple Inc. \textit{Accelerated PyTorch training on Mac}. Apple Developer, 2024. Available at: \url{https://developer.apple.com/metal/pytorch/}

\bibitem{airformat}
SamoZ256. \textit{Breaking down Metal's intermediate representation format}. Medium, 2024. Available at: \url{https://medium.com/@samuliak/breaking-down-metals-intermediate-representation-format-41827022489c}

\bibitem{metaltools}
Apple Inc. \textit{Metal Programming Guide}. Apple Developer Documentation Archive, 2024. Available at: \url{https://developer.apple.com/library/archive/documentation/Miscellaneous/Conceptual/MetalProgrammingGuide/}
\end{thebibliography}

\end{document}
